\documentclass[12pt,a4paper]{article}

\usepackage[T1]{fontenc}
\usepackage[utf8]{inputenc}
% \usepackage[brazil]{babel}  % comentado: babel-portugues não disponível neste TeX Live
\usepackage{amsmath}
\usepackage{amssymb}
\usepackage{listings}
\usepackage{xcolor}
\usepackage[hidelinks]{hyperref}
\usepackage{geometry}

\geometry{margin=2.5cm}

% ----- Estilo para listagens de código Rust -----
\definecolor{codebg}{rgb}{0.97,0.97,0.97}
\definecolor{codekw}{rgb}{0.40,0.20,0.60}
\definecolor{codestr}{rgb}{0.30,0.55,0.30}
\definecolor{codecmt}{rgb}{0.50,0.50,0.50}
\definecolor{codetype}{rgb}{0.20,0.45,0.65}

\lstdefinelanguage{Rust}{
  morekeywords={
    let,mut,fn,pub,use,mod,struct,enum,impl,trait,for,while,loop,if,else,match,
    return,break,continue,as,in,ref,move,static,const,where,self,Self,
    crate,super,box,unsafe,extern,type,dyn
  },
  morecomment=[l]{//},
  morecomment=[s]{/*}{*/},
  morestring=[b]{"}
}

\lstset{
  language=Rust,
  basicstyle=\ttfamily\footnotesize,
  keywordstyle=\color{codekw}\bfseries,
  commentstyle=\color{codecmt}\itshape,
  stringstyle=\color{codestr},
  backgroundcolor=\color{codebg},
  frame=single,
  rulecolor=\color{black!20},
  numbers=left,
  numberstyle=\tiny\color{black!50},
  breaklines=true,
  showstringspaces=false,
  tabsize=2,
  captionpos=b,
  xleftmargin=1em,
  upquote=true,
  literate={á}{{\'a}}1 {ã}{{\~a}}1 {â}{{\^a}}1
           {é}{{\'e}}1 {ê}{{\^e}}1
           {í}{{\'i}}1
           {ó}{{\'o}}1 {ô}{{\^o}}1 {õ}{{\~o}}1
           {ú}{{\'u}}1
           {ç}{{\c{c}}}1
           {Á}{{\'A}}1 {É}{{\'E}}1 {Í}{{\'I}}1 {Ó}{{\'O}}1 {Ú}{{\'U}}1
}

% ----- Metadados -----
\title{\Large\textbf{Relatório Técnico} \\[0.4em]
  \large Trabalho 1 -- DIM0164 Compiladores 2026.1 \\
  Pré-processador, Analisador Léxico e Analisador Sintático para
  uma variante de Java}
\author{João Pedro Araújo, João Paulo Oliveira, Gabriel Sebastião}
\date{\today}

\begin{document}
\maketitle

\begin{abstract}
\noindent Este relatório descreve a implementação em Rust de um \emph{frontend}
de compilador composto por três fases: pré-processador, analisador léxico e
analisador sintático. A linguagem-fonte é uma variante simplificada de Java
(MiniJava) definida na disciplina. Documenta-se a fundamentação teórica das
fases, as estruturas de dados escolhidas, as decisões de design, as
transformações aplicadas à gramática e os mecanismos de reporte de erros.
\end{abstract}

\section{Introdução: Fundamentação Teórica}

O projeto de um compilador é tradicionalmente dividido em duas grandes
partes: \textbf{análise} e \textbf{síntese}. A fase de análise mapeia o
código-fonte para representações gramaticais bem definidas, gera um código
intermediário e indica erros ao programador por meio de mensagens
específicas. A fase de síntese transforma o código intermediário em código
final e, opcionalmente, realiza otimizações~\cite{aho2006compilers}.

A análise propriamente dita é, por sua vez, escalonada em estágios:

\begin{enumerate}
  \item \textbf{Pré-processamento} (\emph{scanning} no Dragon Book) --
    elimina comentários e normaliza espaços em branco. Cuida do código
    bruto antes que ele alcance os autômatos do analisador léxico.
  \item \textbf{Análise léxica} -- agrupa caracteres em \emph{lexemas},
  popula a \emph{tabela de símbolos} para fases posteriores e produz
  \emph{tokens} no formato $\langle\text{text-name}, \text{attribute-value}\rangle$~\cite[\S3.1.3]{aho2006compilers}.
  \item \textbf{Análise sintática} -- combina os tokens em uma árvore
    sintática (explícita ou implícita) que respeita a gramática
    livre-de-contexto da linguagem.
  \item \textbf{Análise semântica} -- verifica propriedades que escapam
    à gramática (tipos, escopos) usando a tabela de símbolos enriquecida.
  \item \textbf{Geração de código intermediário e final} --
    propósito da fase de síntese.
\end{enumerate}

Este trabalho cobre os três primeiros estágios. As fases posteriores ficam
para trabalhos subsequentes da disciplina.

\paragraph{Por que separar pré-processamento e análise léxica?}
Embora os dois manipulem o fluxo bruto de caracteres, separar
responsabilidades torna o código mais legível e localiza decisões
ortogonais. O pré-processador é um \emph{transdutor} simples: lê uma
\texttt{String} e produz outra. O analisador léxico opera sobre essa
\texttt{String} já saneada e produz uma estrutura tipada (lista de
\texttt{Token}). A pequena duplicação de lógica de classificação de
caracteres compensa pela clareza arquitetural.

\paragraph{Por que regex em vez de DFA explícito?}
O Dragon Book apresenta dois caminhos para análise léxica: construir um
autômato finito determinístico (DFA) manualmente, ou descrever os padrões
com expressões regulares e deixar um gerador (\texttt{lex}, \texttt{flex})
transformá-las em DFA. Optei por um meio-termo: usar a crate
\texttt{regex} de Rust diretamente para verificar prefixos contra cada
classe de token, sem gerar um DFA explícito. A motivação foi
pragmática --- tempo curto, recomendação do professor e foco
não-prioritário em construção manual de autômatos.

\section{Visão Geral da Arquitetura}

O \emph{frontend} é organizado em três módulos Rust, mais um \emph{driver}
em \texttt{main.rs}:

\begin{center}
\begin{tabular}{l|l}
\textbf{Módulo} & \textbf{Responsabilidade} \\ \hline
\texttt{preprocessor} & Remove comentários e whitespace redundante \\
\texttt{lexical\_analyzer} & Produz \texttt{Vec<Token>} e \texttt{SymbolTable} \\
\texttt{syntatic\_analyzer} & Verifica corretude sintática \\
\texttt{main} & Compõe a \emph{pipeline} e formata a saída \\
\end{tabular}
\end{center}

A composição vive em \texttt{src/main.rs}:

\begin{lstlisting}[caption={Pipeline em \texttt{main.rs}}]
fn compile(source: &str) -> Result<SymbolTable, CompileError> {
    let preprocessed = preprocessor::preprocess(source)?;
    let (tokens, symbol_table) = lexical_analyzer::get_tokens(&preprocessed);
    syntatic_analyzer::parse(&tokens)?;
    Ok(symbol_table)
}
\end{lstlisting}

\section{Pré-processador}

O pré-processador cumpre os requisitos da spec: remove comentários de
linha (\verb|//|) e de bloco (\verb|/* */|), além de normalizar
whitespace em excesso. Devolve uma \texttt{String} (não um arquivo)
acompanhada de um \textbf{SourceMap}.

\subsection{Decisão: Iteração manual e \texttt{SourceMap}}

Inicialmente, o pré-processador utilizava expressões regulares para
substituição de texto. Embora eficiente para a transformação, essa
abordagem dificultava a rastreabilidade da coluna original, pois a
\texttt{regex} não expunha facilmente o mapeamento entre o deslocamento
do texto resultante e as coordenadas de origem.

Para garantir a fidelidade de mensagens de erro e da tabela de
símbolos, migrei para uma varredura manual caractere a caractere que
constrói um \texttt{SourceMap}:

\begin{lstlisting}[caption={Estrutura de mapeamento de origem}]
pub struct SourceMap {
    pub mapping: Vec<(usize, usize)>,
}

impl SourceMap {
    pub fn get(&self, offset: usize) -> (usize, usize) {
        if offset < self.mapping.len() {
            self.mapping[offset]
        } else {
            self.mapping.last().copied().unwrap_or((1, 1))
        }
    }
}
\end{lstlisting}

A cada caractere preservado ou espaço inserido, registramos sua linha
e coluna originais. Isso permite que fases posteriores (léxica e
sintática) reportem posições na fonte original sem conhecer os detalhes
da minificação.

\subsection{Classes W, P, S e a regra da não-aglutinação}

Para que o pré-processador não funda lexemas adjacentes, ele precisa de
um nível mínimo de consciência léxica. Adoto a divisão clássica em três
classes de \emph{símbolos atômicos}:

\begin{itemize}
  \item \textbf{W} (\emph{word}): caracteres alfanuméricos e \verb|_|.
  \item \textbf{P} (\emph{punctuation}): operadores e delimitadores.
  \item \textbf{S} (\emph{space}): espaços, tabs, quebras de linha.
\end{itemize}

A armadilha está em \texttt{W+S+W}: o espaço entre dois caracteres da
classe W \emph{deve ser preservado}, do contrário palavras-chave se
fundem com identificadores (\verb|int x| viraria \verb|intx|). Para os
demais casos (\texttt{W+S+P}, \texttt{P+S+W}, \texttt{P+S+P}) a remoção
é segura porque a fronteira já é discernível pelo próprio P.

\begin{lstlisting}[caption={Normalização de whitespace com SourceMap}]
pub fn remove_unnecessary_white_spaces(src: &str, incoming_map: SourceMap) -> (String, SourceMap) {
    // ...
    if prev_word && next_word {
        output.push(' ');
        mapping.push(incoming_map.get(byte_offset));
    }
    // ...
}
\end{lstlisting}

\subsection{Remoção de Comentários e Preservação de Posição}

Um comentário de bloco multi-linha é substituído por suas quebras de
linha internas para manter a sincronia de linhas. A grande diferença da
nova implementação é que, para cada caractere preservado (incluindo as
quebras de linha), mapeamos a posição exata (linha, coluna) que ele
ocupava no arquivo original.

\begin{lstlisting}[caption={Rastreamento em comentários de bloco}]
while let Some((_, inner_c)) = char_indices.next() {
    if inner_c == '\n' {
        output.push('\n');
        mapping.push((line, current_column));
        line += 1;
        current_column = 1;
    } else if inner_c == '*' && peek == Some('/') {
        // ...
    }
}
\end{lstlisting}

\subsection{Detecção de \texttt{/*} sem fechamento}

O diagnóstico de comentário aberto agora é feito durante a própria
varredura manual. Se o fim do arquivo é atingido enquanto o estado
indica estar dentro de um comentário de bloco, o erro é emitido com a
linha inicial do comentário (que foi salva no início da detecção do
marcador).

\section{Analisador Léxico}

\subsection{Estruturas de Dados}

A estrutura central segue a convenção do Dragon Book
(\S3.1.3, Tabela 3.4) com um \emph{token name} (classe abstrata) e um
\emph{attribute value} (ponteiro para a tabela de símbolos ou o lexema
embutido):

\begin{lstlisting}[caption={\texttt{Token}, \texttt{TokenClass} e \texttt{TokenAttribute}}]
#[derive(Debug, Clone, PartialEq)]
pub enum TokenClass {
    ID,
    NUMBER,
    KEYWORD(String),
    OPERATION,
    DELIMITER,
    EOF,
    UNKNOWN,
}

#[derive(Debug, Clone, PartialEq)]
pub enum TokenAttribute {
    Pointer(usize),
    Itself(String),
    Null,
}

#[derive(Debug, Clone)]
pub struct Token {
    pub token_name: TokenClass,
    pub attribute_value: TokenAttribute,
    pub line: usize,
    pub column: usize,
}
\end{lstlisting}

A escolha de carregar \texttt{line} e \texttt{column} \emph{em cada
token} é importante: o analisador sintático produz mensagens de erro
com posição exata sem precisar consultar a fonte original.

\subsection{Tabela de Símbolos com \emph{Interning}}

A tabela de símbolos combina dois containers: um \texttt{Vec} para
acessos por índice (o \emph{ponteiro} embutido nos tokens) e um
\texttt{HashMap} para detectar reaparecimentos de um mesmo lexema em
tempo $O(1)$ amortizado. Esse padrão é conhecido como \emph{string
interning} e é o mesmo usado por compiladores reais (LLVM, rustc).

\begin{lstlisting}[caption={\texttt{SymbolTable}, \texttt{SymbolTableEntry} e \texttt{intern()}}]
#[derive(Debug, Clone, PartialEq)]
pub enum SymbolKind {
    Variable,
    Number,
}

#[derive(Debug, Clone)]
pub struct SymbolTableEntry {
    pub lexeme: String,
    pub kind: SymbolKind,
    pub type_info: Option<String>,
    pub first_line: usize,
    pub first_column: usize,
}

#[derive(Debug)]
pub struct SymbolTable {
    pub registers: Vec<SymbolTableEntry>,
    intern_map: HashMap<String, usize>,
}

impl SymbolTable {
    fn intern(&mut self, lexeme: &str, kind: SymbolKind,
              line: usize, column: usize) -> usize {
        if let Some(&idx) = self.intern_map.get(lexeme) {
            return idx;
        }
        let idx = self.registers.len();
        self.registers.push(SymbolTableEntry {
            lexeme: lexeme.to_string(),
            kind,
            type_info: None,
            first_line: line,
            first_column: column,
        });
        self.intern_map.insert(lexeme.to_string(), idx);
        idx
    }
}
\end{lstlisting}

O campo \texttt{type\_info} fica reservado para o analisador semântico
preencher em fase posterior --- o léxico não tem informação suficiente
para inferir tipos.

\subsection{Algoritmo de Longest-Match}

Para cada token, o léxico tenta estender o lexema enquanto houver
\emph{alguma} classe de token cuja regex (ancorada) case com o lexema
completo. Quando estender deixa de produzir match, faz-se \emph{back-up}
de um caractere e o token é emitido com a última classe válida:

\begin{lstlisting}[caption={Núcleo de \texttt{get\_token} com \texttt{SourceMap}}]
fn get_token(full_src: &str, cursor: usize, map: Option<&SourceMap>,
             symbol_table: &mut SymbolTable, line: &mut usize,
             column: &mut usize) -> Result<(Option<Token>, usize), ...> {
    // ...
    if let Some(m) = map {
        let pos = m.get(cursor);
        *line = pos.0; *column = pos.1;
    }
    let mut start_line = *line;
    let mut start_column = *column;

    for c in src.chars() {
        // ...
        if let Some(m) = map {
            let pos = m.get(cursor + consumed);
            *line = pos.0; *column = pos.1;
        } else {
            *column += 1;
        }
        // ... match regex ...
    }
}
\end{lstlisting}

\textbf{Anchorage e disambiguação.} As regexes usam \verb|^...$| para
exigir match completo. Isso resolve a confusão clássica entre
identificador e palavra-chave: \texttt{class1} casa só com a regex
de identificador (não com a de keyword), \texttt{class} casa com as
duas --- a ordem do cascade (\texttt{KEYWORD} antes de \texttt{ID})
privilegia a palavra-chave.

\textbf{Caso especial para \texttt{\&\&}.} Como \verb|&| sozinho não
casa nenhuma classe ancorada, um token de dois caracteres como
\verb|&&| quebraria o algoritmo. A solução pragmática é detectar o
prefixo de um caractere e estender:
\verb|else if lexeme == "&" { continue; }|. É \emph{brittle} para
gramáticas com mais operadores multi-caractere, mas funciona para a
nossa, em que \verb|&&| é o único caso.

\textbf{Contador \texttt{consumed}.} Whitespace e newlines descartados
\emph{antes} do início do lexema ainda foram lidos pelo cursor. Manter
um contador separado de bytes consumidos (que inclui os descartes)
evita off-by-one ao reportar o avanço para o \emph{caller}.

\subsection{Tokens Especiais e Reportagem de Erros}

\paragraph{EOF.} Após o laço principal, \texttt{get\_tokens} sempre
empurra um token \texttt{TokenClass::EOF}. O parser pode então olhar
sempre o próximo token sem checar bound.

\paragraph{Sugestões para identificadores inválidos.} Quando um lexema
não casa nenhuma classe, calculamos sua distância de Levenshtein
contra a lista de palavras-chave (via crate \texttt{strsim}) e, se
houver alguma a no máximo $2$ de distância, incluímos a sugestão na
mensagem:

\begin{lstlisting}[caption={Sugestão por distância de edição}]
fn suggest_keyword(lexeme: &str) -> Option<String> {
    if lexeme.is_empty() { return None; }
    KEYWORDS_LIST
        .iter()
        .map(|kw| (*kw, strsim::levenshtein(lexeme, kw)))
        .filter(|(_, d)| *d <= 2)
        .min_by_key(|(_, d)| *d)
        .map(|(kw, _)| kw.to_string())
}
\end{lstlisting}

Exemplo: dado \verb|claas|, o usuário vê
\texttt{Unknown lexeme: 'claas'. Did you mean: 'class'?}.

\section{Analisador Sintático}

\subsection{Transformação da Gramática}

A spec exige que a gramática seja preparada por \emph{eliminação de
recursão à esquerda} e \emph{fatoração à esquerda}. O documento
\texttt{specs/gramatica-transformada.md} registra todas as alterações;
três pontos merecem destaque:

\paragraph{1. Recursão à esquerda em \texttt{Exp}.} O esquema padrão
$A \to A\alpha \mid \beta$ vira $A \to \beta A'$, $A' \to \alpha A'
\mid \lambda$. A produção original

\begin{verbatim}
Exp -> Exp '&&' Exp | Exp '>' Exp | Exp '+' Exp | ... | Id | Number
\end{verbatim}

\noindent fica:

\begin{verbatim}
Exp     -> ExpBase ExpRest
ExpBase -> 'new' NewRest | '!' Exp | '(' Exp ')'
         | 'true' | 'false' | 'this' | Id | Number
ExpRest -> '&&' Exp ExpRest | '<' Exp ExpRest | '>' Exp ExpRest
         | '+' Exp ExpRest | '-' Exp ExpRest | '*' Exp ExpRest
         | '[' Exp ']' ExpRest | '.' DotRest ExpRest | $\lambda$
\end{verbatim}

\paragraph{2. Sequência de comandos.} A gramática original admite
apenas \emph{um único} \texttt{Cmd} por bloco. Como os arquivos de
teste (\texttt{prog-bubblesort.ling}, \texttt{prog-factorial.ling})
exigem sequências, introduzi:

\begin{verbatim}
Cmds -> Cmd Cmds | $\lambda$
\end{verbatim}

\noindent e substituí \texttt{Cmd} por \texttt{Cmds} no corpo de
\texttt{MainC}, em \texttt{DefMet} e dentro de \texttt{Cmd -> '\{' Cmds '\}'}.

\paragraph{3. Fatorações pontuais.} Em \texttt{Type} entre
\texttt{int} e \texttt{int []}; em \texttt{Cmd} no prefixo
\texttt{Id} (separando \texttt{=} de \texttt{[ Exp ]}); em
\texttt{ExpBase} no prefixo \texttt{new}; em \texttt{ListExp} com a
alternativa $\lambda$ promovida ao chamador.

\paragraph{4. Operador \texttt{<}.} A \texttt{gramatica.md} original
lista apenas \verb|>| como operador relacional, mas
\texttt{bubblesort.ling} e \texttt{factorial.ling} usam \verb|<|.
Acrescentei \verb|<| simetricamente ao \verb|>|.

\subsection{Estratégia: Descendente Recursivo}

Optei por um \textbf{analisador descendente recursivo}. As alternativas
seriam um \emph{table-driven} preditivo (LL(1) com tabela explícita) ou
um gerador (ex.: \texttt{lalrpop}). O descendente recursivo escolhi
por três razões:

\begin{itemize}
  \item \textbf{Legibilidade}: cada não-terminal vira uma função;
    o fluxo do código espelha a gramática.
  \item \textbf{Reportagem de erros}: mais flexível que tabela
    preditiva pura --- dá pra fabricar mensagens contextuais.
  \item \textbf{Sem dependências}: gerador adicionaria peso ao
    \emph{build}.
\end{itemize}

\subsection{Estrutura}

O parser carrega apenas o slice de tokens e um cursor:

\begin{lstlisting}[caption={Estrutura do parser}]
pub struct Parser<'a> {
    tokens: &'a [Token],
    cursor: usize,
}

#[derive(Debug, PartialEq)]
pub struct ParseError {
    pub line: usize,
    pub column: usize,
    pub msg: String,
}
\end{lstlisting}

Cada não-terminal da gramática transformada corresponde a um método
privado. Em \texttt{Cmd} a estrutura é especialmente clara:

\begin{lstlisting}[caption={Trecho de \texttt{parse\_cmd}}]
fn parse_cmd(&mut self) -> Result<(), ParseError> {
    let t = self.peek();
    if Self::is_keyword(t, "if") {
        self.advance();
        self.expect_delim("(")?;
        self.parse_exp()?;
        self.expect_delim(")")?;
        self.parse_cmd()?;
        self.expect_keyword("else")?;
        self.parse_cmd()?;
        return Ok(());
    }
    if Self::is_keyword(t, "while") {
        self.advance();
        self.expect_delim("(")?;
        self.parse_exp()?;
        self.expect_delim(")")?;
        self.parse_cmd()?;
        return Ok(());
    }
    // ... outros ramos ...
}
\end{lstlisting}

\subsection{Disambiguação por \emph{Lookahead}}

Duas situações exigem mais de um token de \emph{lookahead}:

\paragraph{DefVar vs Cmd com \texttt{Id} no início.} A gramática
permite \texttt{Type Id ;} (declaração com tipo custom) ou
\texttt{Id = Exp ;} (atribuição). Ambos começam com um \texttt{Id}.
A regra usa o token seguinte: se for outro \texttt{Id}, é
declaração; se for \texttt{=} ou \texttt{[}, é comando.

\begin{lstlisting}[caption={Discriminação por \emph{lookahead}-2}]
fn at_def_var(&self) -> bool {
    let t = self.peek();
    if Self::is_keyword(t, "int") || Self::is_keyword(t, "boolean") {
        return true;
    }
    if Self::is_id(t) {
        return Self::is_id(self.peek_at(1));
    }
    false
}
\end{lstlisting}

\paragraph{\texttt{DotRest}: \texttt{.length} vs \texttt{.metodo(args)}.}
Após \verb|.| dentro de uma expressão, o próximo token pode ser a
palavra-chave \texttt{length} (acesso) ou um \texttt{Id} seguido de
\verb|(...)| (chamada de método). \emph{Lookahead}-1 sobre o tipo do
próximo token resolve.

\subsection{Reportagem de Erros}

A função \texttt{error} captura o token corrente e o estiliza:

\begin{lstlisting}[caption={Erro com posição e descrição do token corrente}]
fn error(&self, msg: String) -> ParseError {
    let tok = self.peek();
    let got = describe_token(tok);
    ParseError {
        line: tok.line,
        column: tok.column,
        msg: format!("{msg}, got {got}"),
    }
}
\end{lstlisting}

\noindent o que dá mensagens como
\verb|expected ';', got '}'| acompanhadas de linha e coluna do
token ofensor. Isso atende ao requisito da spec de informar
``exatamente o local e o tipo do erro''.

\subsection{Precedência de Operadores}

Em uma transformação ingênua, o esquema $A \to A\alpha \mid \beta$
elimina recursão à esquerda mas perde precedência: \texttt{a + b * c}
seria parseado como \texttt{(a+b)*c}. Para preservar a precedência
matemática (e o que é natural pra um programador de C/Java), a
gramática de expressões foi \textbf{estratificada em níveis}, do
menor para o maior:

\begin{verbatim}
Exp        -> ExpAnd
ExpAnd     -> ExpRel  ( '&&' ExpRel )*
ExpRel     -> ExpAdd  ( ('<'|'>') ExpAdd )*
ExpAdd     -> ExpMul  ( ('+'|'-') ExpMul )*
ExpMul     -> ExpUnary( '*'      ExpUnary )*
ExpUnary   -> '!' ExpUnary | ExpPostfix
ExpPostfix -> ExpPrimary ( '[' Exp ']' | '.' DotRest )*
ExpPrimary -> 'new' NewRest | '(' Exp ')'
            | 'true' | 'false' | 'this' | Id | Number
\end{verbatim}

Cada nível associa à esquerda via \emph{loop} iterativo (não recursão
à direita) e delega ao nível imediatamente superior em precedência.

\begin{lstlisting}[caption={Nível \texttt{ExpAdd}: \texttt{+} e \texttt{-} associando à esquerda}]
    fn parse_exp_add(&mut self) -> Result<(), ParseError> {
        self.parse_exp_mul()?;
        while Self::is_op_str(self.peek(), "+") || Self::is_op_str(self.peek(), "-") {
            self.advance();
            self.parse_exp_mul()?;
        }
        Ok(())
    }
}
\end{lstlisting}

\section{Resultados}

\subsection{Cobertura de Testes}

A suíte de testes de integração tem \textbf{136 testes verde}
distribuídos em quatro arquivos:

\begin{center}
\begin{tabular}{l|r}
\textbf{Suíte} & \textbf{Testes} \\ \hline
\texttt{tests/preprocessor.rs} & 37 \\
\texttt{tests/lexical\_analyzer.rs} & 52 \\
\texttt{tests/syntatic\_analyzer.rs} & 38 \\
\texttt{tests/source\_position.rs} & 9 \\
\hline
\textbf{Total} & \textbf{136} \\
\end{tabular}
\end{center}

\noindent A suíte \texttt{tests/source\_position.rs} é especialmente
importante: ela garante que, mesmo após a minificação agressiva do
pré-processador, os tokens retêm suas coordenadas originais na fonte.
A suíte do parser inclui testes específicos verificando o
...\textbf{shape da AST} (precedência aritmética, associatividade
esquerda de \texttt{+/-}, \texttt{*} ligando mais forte que
\texttt{+}, \texttt{\&\&} ligando mais fraco que \texttt{</>}, e
parênteses sobrepondo precedência).

\noindent As suítes do léxico cobrem cada classe de token, casos de
\emph{longest-match}, sequências multi-token, \emph{interning} de
identificadores e rastreamento de linha/coluna. As do parser cobrem
programas mínimos, todos os tipos de comandos, todas as formas de
expressão, casos de erro com verificação de linha/coluna e
\emph{end-to-end} dos fixtures \texttt{bubblesort.ling} e
\texttt{factorial.ling}.

\subsection{Programas de Exemplo}

Os dois programas exemplos da disciplina passam pelo pipeline
completo:

\begin{verbatim}
$ cargo run -- specs/prog-bubblesort.ling
código está sintaticamente correto

tabela de símbolos:
     #  lexeme                    kind        type        line     col
     0  BubbleSort                Variable  -              1       7
     ...

$ cargo run -- specs/prog-factorial.ling
código está sintaticamente correto

tabela de símbolos:
     #  lexeme                    kind        type        line     col
     0  Factorial                 Variable  -              1       7
     1  a                         Variable  -              2      33
     ...
\end{verbatim}

\section{Conclusão}

O \emph{frontend} resultante atende aos requisitos do trabalho:

\begin{itemize}
  \item Pré-processador remove comentários e whitespace excessivo,
    preserva quebras de linha de comentários para rastreamento
    posterior, detecta \verb|/*| não fechado.
  \item Analisador léxico produz tokens com classe, atributo, linha e
    coluna, mantém uma tabela de símbolos com \emph{interning},
    sugere correções para identificadores tipados errados.
  \item Analisador sintático descendente recursivo opera sobre a
    gramática transformada (recursão à esquerda removida, fatoração à
    esquerda aplicada), reporta erros com localização exata.
\end{itemize}

A escolha de Rust contribuiu para o resultado em duas direções: o
sistema de tipos forçou clareza nas fronteiras entre fases
(\texttt{Result} obriga a tratar erros) e a crate \texttt{regex} (com
suporte nativo a classes POSIX e operações entre classes) reduziu
significativamente a quantidade de código manual de varredura.

Como trabalhos futuros, ficam pendentes: a fase de análise semântica
(preenchimento do campo \texttt{type\_info} da tabela de símbolos com
tipos inferidos), a construção explícita de AST se a fase seguinte
exigir, e a geração de código intermediário.

\begin{thebibliography}{9}
\bibitem{aho2006compilers}
A.~V. Aho, M.~S. Lam, R.~Sethi, J.~D. Ullman.
\emph{Compilers: Principles, Techniques, and Tools (Dragon Book)}.
2.~ed. Addison-Wesley, 2006.
\end{thebibliography}

\end{document}
